\begin{center}\begin{tikzpicture}[font=\small,>=stealth,line join=round]
\draw[gray!40,rounded corners=3pt] (0,0) rectangle (14.2,8.4);\node at (7.1,7.800000000000001) {证明一：区间不动，先在数轴上找最近与最远};
\begin{scope}[shift={(2.5,6.1)},x=1.15cm]\draw[->,gray] (-1.6,0)--(5.8,0);\draw[blue,line width=2pt] (0,0)--(4,0);\foreach \v in {0,4}{\fill[blue] (\v,0) circle(2pt);\node[below=5pt] at (\v,0){$\v$};}\fill[violet] (-1,0) circle(2.5pt);\node[above=2pt] at (-1,0){$h=-1$};\draw[orange!85!black,line width=2pt] (0,-.16)--(0,.16);\node[anchor=west] at (-1.6,.95){轴在左侧：最近$0$；最远4};\end{scope}\begin{scope}[shift={(2.5,4.1)},x=1.15cm]\draw[->,gray] (-1.6,0)--(5.8,0);\draw[blue,line width=2pt] (0,0)--(4,0);\foreach \v in {0,4}{\fill[blue] (\v,0) circle(2pt);\node[below=5pt] at (\v,0){$\v$};}\fill[violet] (2,0) circle(2.5pt);\node[above=2pt] at (2,0){$h=2$};\draw[orange!85!black,line width=2pt] (2,-.16)--(2,.16);\node[anchor=west] at (-1.6,.95){轴在区间内：最近$2$；最远0与4};\end{scope}\begin{scope}[shift={(2.5,2.1)},x=1.15cm]\draw[->,gray] (-1.6,0)--(5.8,0);\draw[blue,line width=2pt] (0,0)--(4,0);\foreach \v in {0,4}{\fill[blue] (\v,0) circle(2pt);\node[below=5pt] at (\v,0){$\v$};}\fill[violet] (5,0) circle(2.5pt);\node[above=2pt] at (5,0){$h=5$};\draw[orange!85!black,line width=2pt] (4,-.16)--(4,.16);\node[anchor=west] at (-1.6,.95){轴在右侧：最近$4$；最远0};\end{scope}
\node[text width=12.6cm,align=center] at (7.1,.65) {蓝段为[0,4]，橙色短标记是最近位置；轴在区间内时可以直接取轴的位置。};
\end{tikzpicture}\end{center}
